\renewcommand{\arraystretch}{1.25}
\captionsetup[longtable]{justification=raggedright,singlelinecheck=false}

\begin{longtable}{|p{1.2cm}|p{3cm}|p{3.8cm}|p{3.8cm}|}
\caption{Matriks Analisis Risiko dan Strategi Mitigasi}
\label{tbl:risk_matrix} \\ \hline
\textbf{Kode} & \textbf{Risiko} & \textbf{Dampak Potensial} & \textbf{Strategi Mitigasi} \\ \hline
\endfirsthead

\multicolumn{4}{l}{Tabel \ref{tbl:risk_matrix} Matriks Analisis Risiko dan Strategi Mitigasi (lanjutan)} \\ \hline
\textbf{Kode} & \textbf{Risiko} & \textbf{Dampak Potensial} & \textbf{Strategi Mitigasi} \\ \hline
\endhead

\hline
\multicolumn{4}{r}{\textit{bersambung ke halaman berikutnya}} \\
\endfoot

\hline
\endlastfoot

\textbf{R-01} &
Ketidaklengkapan Data Swagger &
Struktur definisi objek tidak sepenuhnya terdeteksi sehingga relasi graf tidak komplet dan \textit{retrieval} menjadi kurang akurat. &
Melakukan \textit{cross-check} manual terhadap objek paling kritikal (Deployment, Pod, Service, ConfigMap), serta menambahkan validasi referensial untuk mendeteksi anomali struktur. \\ \hline

\textbf{R-02} &
Ambiguitas Relasi Graf &
\textit{Edge} tertentu tidak terbaca pada pola kompleks (\textit{nested, selector, injection}) sehingga traversal kehilangan jalur valid dan menurunkan skor \textit{precision/recall}. &
Menambahkan aturan \textit{fallback traversal}, memperkaya heuristik pengenalan pola relasi, dan melakukan uji representatif terhadap sampel graf untuk memastikan ketepatan pemodelan secara keseluruhan. \\ \hline

\textbf{R-03} &
Kegagalan Generasi Kode (\textit{Syntax Error}) &
Konfigurasi yang dihasilkan tidak dapat diparse dan berpotensi gagal dieksekusi, sehingga melanggar standar \textit{Syntactic Validity}. &
Mengintegrasikan PyYAML sebagai validator otomatis, serta menambahkan prosedur \textit{post-processing} sebagai prosedur koreksi jika parser mendeteksi kesalahan sintaksis. \\ \hline

\textbf{R-04} &
Bias Penilaian oleh LLM \textit{Judge} &
Evaluasi performa dapat bias jika LLM penilai menghasilkan keputusan tidak konsisten atau tidak faktual. &
Menggunakan model penilai yang lebih kuat, melakukan \textit{cross-validation} sebagian sampel secara manual, serta membandingkan referensi dengan dokumentasi resmi Kubernetes. \\ \hline

\textbf{R-05} &
Latensi Tinggi pada Proses \textit{Retrieval--Generation} &
Waktu respon melebihi batas toleransi (maksimal 60 detik), sehingga sistem sulit digunakan secara \textit{real-time}. &
Mengoptimalkan \textit{pipeline} dengan \textit{caching}, pembatasan traversal, \textit{batching request}, serta melakukan stress test untuk memastikan waktu rata-rata berada dalam batas KPI. \\ \hline

\textbf{R-06} &
\textit{Hallucination} pada Mode \textit{Code Generation} &
LLM menghasilkan YAML yang tidak terdapat dalam skema Kubernetes, menurunkan \textit{faithfulness} dan membahayakan konfigurasi. &
Menambahkan \textit{negative constraints} dalam \textit{prompt}, membatasi penggunaan field hanya dari \textit{Knowledge Graph}, dan menerapkan mekanisme \textit{fail-safe} jika ditemukan properti asing. \\ \hline

\end{longtable}